% LuaLaTeXで日本語文書を作成するためのドキュメントクラス (Document Class) を指定します。
\documentclass[12pt]{ltjsarticle}

% LuaTeX-jaにおいて、フォントの指定や切り替えを行うためのパッケージ (Package) を読み込みます。
\usepackage{luatexja-fontspec}

% geometry パッケージを使って、用紙サイズと余白を設定します
\usepackage[
  a4paper,      % Paper Size（用紙サイズ）: a4paper, b5paper, letterpaper など
  top=25mm,     % Top Margin（上余白）
  bottom=25mm,  % Bottom Margin（下余白）
  left=20mm,    % Left Margin（左余白）
  right=20mm    % Right Margin（右余白）
]{geometry}

% --------------------------------------------------
% メインフォントの候補（メモ書き）
% --------------------------------------------------
%IPAmjMincho
%Harano Aji Mincho
%Source Han Serif

% --------------------------------------------------
% メインフォントの設定
% --------------------------------------------------
% 1. メインフォントの宣言
% 文書全体の基準となる日本語フォント (Main Japanese Font) を
%「原ノ味明朝 (Harano Aji Mincho)」に指定します。
% Language=Japanese を指定することで、日本語向けの適切な字形や組版が適用されます。
\setmainjfont{Harano Aji Mincho}[
  Language=Japanese
]

% --------------------------------------------------
% 拡張漢字用の代替フォント（フォールバック用）の定義
% --------------------------------------------------
% 2. 代替用（Ext.B 以降など）のフォントファミリーを新しく定義する
% 原ノ味明朝に収録されていない非常に珍しい漢字（拡張漢字）を手動で出力するために、
% それぞれ専用のフォントを呼び出すためのコマンド (Command) を定義します。

% 注意: CJK漢字は「和文 (Japanese/CJK)」として扱われるため、
% \newfontfamily（欧文用）ではなく \newjfontfamily（和文用）を使います。

% \HanaBFont というコマンドで「HanaMinB (花園明朝B)」を使えるようにします。主として拡張Bなどの漢字用です。
\newjfontfamily\HanaBFont{HanaMinB}
% \JigmoIIFont というコマンドで「Jigmo2」を使えるようにします。
\newjfontfamily\JigmoIIFont{Jigmo2}
% \JigmoIIIFont というコマンドで「Jigmo3」を使えるようにします。主として拡張Gなどさらに新しい漢字用です。
\newjfontfamily\JigmoIIIFont{Jigmo3}

% --------------------------------------------------
% 文書の本体 (Document Body)
% --------------------------------------------------
\begin{document}

\section*{フォント出力テスト (Font Output Test)}

% --------------------------------------------------
\subsection*{1. Harano Aji Mincho (原ノ味明朝)}
% 通常のテキストは自動的に原ノ味明朝（メインフォント）で出力されます。
あいうえお かきくけこ さしすせそ たちつてと なにぬねの
はひふへほ まみむめも やゆよ らりるれろ わをん
アイウエオ カキクケコ サシスセソ タチツテト ナニヌネノ
ハヒフヘホ マミムメモ ヤユヨ ラリルレロ ワヲン
常用漢字：海、山、川、空、木、林、森、花、草、虫、犬、猫、鳥、魚、貝
教育漢字：右左上下中大小春夏秋冬東西南北年月日水金地火木土天
JIS第1・第2水準漢字：魑魅魍魎、薔薇、憂鬱、躊躇、麒麟、挨拶、完璧
四字熟語：一期一会、温故知新、花鳥風月、喜怒哀楽、明鏡止水、森羅万象
記号類：！？＠＃＄％＆＊（）「」『』【】＋－×÷＝≒
追加記号：〇✖△▽♨♪♫♬♠♡♣♢★☆※〒〠

% --------------------------------------------------
\subsection*{2. HanaMinB (花園明朝B)}
% 主に CJK統合漢字拡張B (Ext.B) などを担当するフォントです。
% フォント切り替えコマンド \HanaBFont を適用します。
{\HanaBFont
𠀋 𠀾 𠁆 𠁕 𠁣 𠂤 𠃓 𠄘 𠄡 𠅘 𠅤 𠆢 𠆤 𠇔 𠇲
𠈓 𠈲 𠉙 𠉛 𠊞 𠊟 𠋥 𠋤 𠌥 𠌤 𠍥 𠍤 𠎥 𠎤 𠏥
𠏤 𠐥 𠐤 𠑥 𠑤 𠒥 𠒤 𠓥 𠓤 𠔥 𠔤 𠕥 𠕤 𠖥 𠖤
𠮟 𡈽 𡌛 𡑮 𡚴 𡢽 𡧃 𡨸 𡫏 𡮣
𡿖 𡿗 𡿘 𡿙 𡿚 𡿛 𡿜 𡿝 𡿞 𡿟 𡿠 𡿡 𡿢 𡿣 𡿤
𢀖 𢀗 𢀘 𢀙 𢀚 𢀛 𢀜 𢀝 𢀞 𢀟 𢀠 𢀡 𢀢 𢀣 𢀤
𢁖 𢁗 𢁘 𢁙 𢁚 𢁛 𢁜 𢁝 𢁞 𢁟 𢁠 𢁡 𢁢 𢁣 𢁤
𢂖 𢂗 𢂘 𢂙 𢂚 𢂛 𢂜 𢂝 𢂞 𢂟 𢂠 𢂡 𢂢 𢂣 𢂤
𢃖 𢃗 𢃘 𢃙 𢃚 𢃛 𢃜 𢃝 𢃞 𢃟 𢃠 𢃡 𢃢 𢃣 𢃤
}

% --------------------------------------------------
\subsection*{3. Jigmo2}
% 主に CJK統合漢字拡張C, D, E, F などの中間的な拡張領域を担当するフォントです。
% フォント切り替えコマンド \JigmoIIFont を適用します。
{\JigmoIIFont
% 拡張Cの文字群
𪜀 𪜁 𪜂 𪜃 𪜄 𪜅 𪜆 𪜇 𪜈 𪜉 𪜊 𪜋 𪜌 𪜍 𪜎
𪜏 𪜐 𪜑 𪜒 𪜓 𪜔 𪜕 𪜖 𪜗 𪜘 𪜙 𪜚 𪜛 𪜜 𪜝
% 拡張Dの文字群
𫝀 𫝁 𫝂 𫝃 𫝄 𫝅 𫝆 𫝇 𫝈 𫝉 𫝊 𫝋 𫝌 𫝍 𫝎
𫝏 𫝐 𫝑 𫝒 𫝓 𫝔 𫝕 𫝖 𫝗 𫝘 𫝙 𫝚 𫝛 𫝜 𫝝
% 拡張Eの文字群
𫠠 𫠡 𫠢 𫠣 𫠤 𫠥 𫠦 𫠧 𫠨 𫠩 𫠪 𫠫 𫠬 𫠭 𫠮
𫠯 𫠰 𫠱 𫠲 𫠳 𫠴 𫠵 𫠶 𫠷 𫠸 𫠹 𫠺 𫠻 𫠼 𫠽
% 拡張Fの文字群
𬺰 𬺱 𬺲 𬺳 𬺴 𬺵 𬺶 𬺷 𬺸 𬺹 𬺺 𬺻 𬺼 𬺽 𬺾
𬺿 𬻀 𬻁 𬻂 𬻃 𬻄 𬻅 𬻆 𬻇 𬻈 𬻉 𬻊 𬻋 𬻌 𬻍
}

% --------------------------------------------------
\subsection*{4. Jigmo3}
% 主に CJK統合漢字拡張G など、さらに新しい・非常に難解な漢字群を担当するフォントです。
% フォント切り替えコマンド \JigmoIIIFont を適用します。
{\JigmoIIIFont
% 拡張Gの文字群
𰀀 𰀁 𰀂 𰀃 𰀄 𰀅 𰀆 𰀇 𰀈 𰀉 𰀊 𰀋 𰀌 𰀍 𰀎
𰀏 𰀐 𰀑 𰀒 𰀓 𰀔 𰀕 𰀖 𰀗 𰀘 𰀙 𰀚 𰀛 𰀜 𰀝
𰀞 𰀟 𰀠 𰀡 𰀢 𰀣 𰀤 𰀥 𰀦 𰀧 𰀨 𰀩 𰀪 𰀫 𰀬
𰀭 𰀮 𰀯 𰀰 𰀱 𰀲 𰀳 𰀴 𰀵 𰀶 𰀷 𰀸 𰀹 𰀺 𰀻
𰻞 (ビィャン) 𰻝 𰻜 𰻛 𰻚 𰻙 𰻘 𰻗 𰻖 𰻕 𰻔 𰻓
𰾏 𰾎 𰾍 𰾌 𰾋 𰾊 𰾉 𰾈 𰾇 𰾆 𰾅 𰾄 𰾃 𰾂 𰾁
𰾐 𰾑 𰾒 𰾓 𰾔 𰾕 𰾖 𰾗 𰾘 𰾙 𰾚 𰾛 𰾜 𰾝 𰾞
𰿀 𰿁 𰿂 𰿃 𰿄 𰿅 𰿆 𰿇 𰿈 𰿉 𰿊 𰿋 𰿌 𰿍 𰿎
}

\end{document}
